\section{Massive and Off-Shell Generalizations}
\label{sec:massive-offshell-bcfw}

\textbf{Literature.} Two directions extend BCFW beyond the massless on-shell setting of Section~\ref{sec:bcfw-recursion}. The \emph{massive} direction rests on the spinor-helicity formalism for arbitrary masses and spins of \autocite{Arkani-Hamed:2017jhn}, which makes the massive little group $SU(2)$ manifest and classifies all three-point amplitudes; massive BCFW-type recursions then follow from combined massless/massive momentum shifts \autocite{Ballav:2020ese}, and from the all-line-transverse and massive-BCFW shifts of \autocite{Ema:2024vww}, applied to the electroweak theory in \autocite{Ema:2024rss}, with applications to supersymmetric \autocite{Abhishek:2022nqv} and Standard-Model \autocite{Christensen:2024bdt} amplitudes, an on-shell mass-insertion reorganization relating massive to massless amplitudes \autocite{Ni:2025xkg}, and a numerical implementation \autocite{Christensen:2024xsg}; a pedagogical introduction covering the massive case is \autocite{bognaScatteringAmplitudes} (the thesis \autocite{colwellDualitiesHelicity2017} reviews the massless formalism, with only a brief massive section). The \emph{off-shell/massless} direction generalizes the shift from the original two-line shift \autocite{Britto:2005fq} to all-line \autocite{Cohen:2010mi} and one-line \autocite{Cohen:2024xuf} shifts, to off-shell external gluons \autocite{vanHameren:2014iua}, and to Wilson-line form factors \autocite{Cohen:2024xuf}. The all-line shift is also related to the MHV vertex expansion proved in \autocite{Elvang:2008vz}. 

\subsection{Massive Spin-Helicity and BCFW}
\label{subsec:massive-spin-helicity-and-bcfw}

The BCFW recursion, after its original formulation, was soon generalized to massive particles with arbitrary spins, in particular fermions. This line of work has started since the original BCFW formulation in 2005, and a recent seminal work is \autocite{Arkani-Hamed:2017jhn}. 
A more recent line of research applies this general spin-helicity formalism to study various types of scattering amplitudes and recursion relations, such as \autocite{Abhishek:2022nqv} and references therein. Also see \autocite{Christensen:2024bdt, Ema:2024rss} and references therein.

\subsection{Massless Off-Shell BCFW}
\label{subsec:massless-off-shell-bcfw}

An early work on the off-shell BCFW is \autocite{vanHameren:2014iua}, for tree-level multi-gluon amplitudes. A recent work with a mathematical review of the BCFW-type construction is \autocite{Cohen:2024xuf}.

In \autocite{Britto:2005fq} two external lines are shifted to complexify the amplitude; there has been generalization to all-line shifts \autocite{Cohen:2010mi}, and one-line shifts \autocite {Cohen:2024xuf}.

\subsection{Massive QED}
\label{subsec:massive-qed}

Massive QED is the simplest testing ground for massive on-shell recursion. \autocite{Ema:2024vww} constructs the four- and five-point massive QED amplitudes from on-shell recursion, showing that the all-line-transverse and massive-BCFW shifts avoid the contact-term ambiguities of a naive ``gluing'' of lower-point amplitudes. \autocite{Ni:2025xkg} isolates the correct ultraviolet of the three-point $F\bar F\gamma$ amplitude through its on-shell mass insertion, and the constructive computation of \autocite{Christensen:2024bdt} handles internal photon lines by giving the photon a mass and taking the massless limit, recovering the ordinary QED results. 